problem:  
Let \((M^n,g)\) be a complete, noncompact Riemannian manifold satisfying **volume doubling** and a **scale-invariant Poincaré inequality**, so that the standard analysis of harmonic functions applies.  
Assume that the Ricci curvature decays at an almost-quadratic rate:
\[
\mathrm{Ric}_g(x) \ge -K(1+r(x))^{-2+\varepsilon}, \quad r(x)=\mathrm{dist}(x,o),
\]
for some constants \(K>0\) and \(\varepsilon>0\).  
Let \(\mathcal{H}_{\mathrm{poly}}(M)\) denote the space of harmonic functions of polynomial growth, and let \(\mathcal{H}_{\mathrm{fin}}(M)\) denote the space of finite-energy harmonic functions.

---

**Open problem:**  
Does any subquadratic improvement in the Ricci curvature decay (\(\varepsilon>0\)) enforce *analytic rigidity* in the vanishing behavior of harmonic functions? More precisely:

1. **Finiteness threshold:**  
   Is there a constant \(\varepsilon_c>0\) depending only on dimension \(n\) and the doubling/Poincaré constants such that if  
   \[
   \mathrm{Ric}_g(x) \ge -K(1+r(x))^{-2+\varepsilon_c},
   \]
   then
   \(\mathcal{H}_{\mathrm{poly}}(M)\) is *finite-dimensional* and  
   \(\mathcal{H}_{\mathrm{fin}}(M)=\mathbb{R}\) (i.e., every finite-energy harmonic function is constant)?

2. **Sharpness and counterexamples:**  
   Conversely, if \(\varepsilon<\varepsilon_c\), can one construct manifolds satisfying the same analytic structure (doubling, Poincaré) and curvature decay but exhibiting **nontrivial** polynomial-growth or finite-energy harmonic functions?

Equivalently, the question asks:  
> Is there a universal **analytic vanishing threshold** just beyond quadratic Ricci curvature decay—analogous to parabolicity thresholds depending on volume growth—such that the harmonic function space collapses to its minimal form precisely when curvature decays slightly faster than \(r^{-2}\)?

---

Why is it a "good" problem:  
This formulation avoids the previous ambiguities by connecting curvature decay directly to analytic structures where classical vanishing theorems operate. It cleanly asks whether *any improvement beyond quadratic Ricci decay* forces geometric–analytic rigidity, bridging two deep themes: curvature decay and Liouville-type phenomena.  
It unifies tools from geometric analysis (Ricci curvature control, volume doubling, and Poincaré inequalities) with potential theory and harmonic growth analysis.  

The problem is simple to state yet conceptually far-reaching: the quadratic decay rate marks a natural geometric boundary where many analytic estimates begin to fail. Discovering a precise threshold \(\varepsilon_c\) would reveal whether harmonic rigidity depends continuously on geometric decay, producing a curvature-driven analogue of the Colding–Minicozzi and Li–Wang dimension estimates.  

It encapsulates the “narrow entrance, profound depth’’ ideal by seeking a fundamental curvature–analysis boundary controlling the existence of harmonic functions—an open frontier in modern geometric analysis.